\subsection{Correlated Telemetry Unavailability (Synthetic Stress Test)}
The previous value-dependent masking study assumes conditionally independent
signal-missingness indicators. We separately stress-test \emph{correlation
between} missingness indicators while holding each signal's expected dropout
probability at 0.25 and making masking independent of signal values. Four
paired regimes range from independent loss to common-mode loss of five
endpoint-related signals and finally all eight signals. The endpoint grouping
and correlation are stipulated, not estimated from enterprise data. Table~\ref{tab:correlated-outage}
reports ordinary-ALLOW rates among scenario-labeled unsafe observations.
\begin{table}[t]
\centering
\caption{Unsafe ordinary ALLOW (\%), synthetic correlated missingness;
three mask seeds on 50,000 fixed generated observations.}
\label{tab:correlated-outage}
\begin{tabular}{lrrrr}
\toprule
Mask regime & Renorm. & Neutral & Pessim. & Step-up \\
\midrule
Independent & 17.47 & 4.54 & 1.46 & 6.07 \\
Endpoint shared 50\% & 20.13 & 5.78 & 2.82 & 7.67 \\
Endpoint shared 100\% & 22.63 & 7.13 & 4.09 & 9.43 \\
All shared 100\% & 9.34 & 9.34 & 9.34 & 9.34 \\
\bottomrule
\end{tabular}
\end{table}
The endpoint-shared 100\% regime increased pessimistic unsafe ALLOW from
1.46\% to 4.09\%, despite similar marginal dropout. Under all-signal
common-mode loss, observations are either complete or wholly missing; all
four policies consequently agree, but 25.125\% of rows have no telemetry
and are escalated. The nonmonotone renormalization response underscores that
marginal missingness rates do not determine decision risk. Safe-user DENY
and STEP\_UP rates are reported separately in the reproducibility artifact.
These generator-defined outcomes do not estimate real-world policy accuracy.
Sensor-network literature describes correlated observations and downtime
\cite{dinwoodie2019sensor}; it does not establish an endpoint outage model.
